\medskip

\textbf{Principaux domaines abordés :} 
probabilités.

\medskip



















On peut compléter l'arbre pondéré :

\begin{minipage}{0.48\linewidth}

\begin{itemize}
\item[$\bullet~~$] $E_0$ : \og le bit envoyé est un $0$ \fg{} ;
\item[$\bullet~~$] $E_1$ : \og le bit envoyé est un 1 \fg{} ;
\item[$\bullet~~$] $R_0$ : \og le bit reçu est un $0$\fg{} 
\item[$\bullet~~$] $R_1$ : \og le bit reçu est un $1$ \fg.
\end{itemize}
\end{minipage}\hfill
\begin{minipage}{0.48\linewidth}
\begin{center}
\pstree[treemode=R,nodesepA=0pt,nodesepB=2.5pt,treesep = 1cm,levelsep=2.5cm]{\TR{}}
{\pstree{\TR{$E_0$~}\taput{0,4}}
{\TR{$R_0$}\taput{\blue 0,99}
\TR{$R_1$}\tbput{0,01}
}
\pstree{\TR{$E_1$~}\tbput{\blue 0,6}}
{\TR{$R_0$}\taput{0,02}
\TR{$R_1$}\tbput{\blue 0,98}
}}
\end{center}
\end{minipage}

On sait que:

$p\left(E_0\right) = 0,4 ;\quad  p_{R_0}\left(R_1\right) = 0,01 ;\quad p_{R_1}\left(R_0\right) = 0,02$.

On rappelle que la probabilité conditionnelle de $A$ sachant $B$ est notée $p_B(A)$.

On peut ainsi représenter la situation par l'arbre de probabilités ci-dessus.

\medskip

\begin{enumerate}
\item 
$p\left(E_0 \cap R_0\right) = p\left(E_0 \right) \times p_{E_0}\left(R_0 \right) = 0,4 \times 0,99 = 0,396$.






\item 
On a aussi $p\left(E_1 \cap R_0\right) = p\left(E_1 \right) \times p_{E_1}\left(R_0 \right) = 0,6 \times 0,02 = 0,012$.

 D'après la loi des probabilités totales : $p\left(R_0\right) = p\left(E_0 \cap R_0\right) + p\left(E_1 \cap R_0\right) = 0,396 + 0,012 = 0,408$.





\item 

Avec $p\left(R_1 \right) = 1 - p\left(R_0\right) = 1 - 0,408 = 0,592$ ;

$p_{R_1}\left(E_0\right) = \dfrac{p\left(R_1 \cap E_0 \right)}{p\left(R_1\right)} = \dfrac{p\left(E_0 \cap R_1 \right)}{p\left(R_1\right)} = \dfrac{0,4 \times 0,01}{0,592} = \dfrac{0,004}{0,592} \approx \np{0,0067}$, soit environ 0,007 au millième près.





\item 





On a $p(\text{erreur de transmission}) = p\left(E_0 \cap R_1\right) + p\left(E_1 \cap R_0\right) = 0,4 \times 0,01 + 0,6 \times 0,02 = 0,004 + 0,012 = 0,016$.
\end{enumerate}

Un message de longueur huit bits est appelé un octet.

On admet que la probabilité qu'un octet soit transmis sans erreur est égale à 0,88.

\begin{enumerate}[resume]
\item 








On a une loi binomiale de paramètres $n = 10$ et $p = 0,88$. Si $x$ est la variable aléatoire égal au nombre d'octets transmis sans erreur, on a :

$p(X = 7) = \binom{10}{7}\times 0,88^7 \times (1 - 0,88)^{10-7} = \binom{10}{7}\times 0,88^7 \times (0,12)^{3} \approx \np{0,0847}$ soit 0,085 au millième près.
\item 








On a $p(X \geqslant 1) = 1 - p(X = 0) = 1 - \binom{10}{0}0,88^{0} \times 0,12^{10} = 1 - 0,12^{10}$.
\item 










On a $p(X = 18) = 0,88^{18} \approx 0,109 > 0,1$. Donc $N_0 = 18$.
\end{enumerate}
